\subsection{Ordering}

\begin{definition}[Greater than]
\label{def:greater-than}
Let \(x,y \in \mathbb{N}\). We write
\[
x > y
\]
if there exists \(u \in \mathbb{N}\) such that
\[
x = y + u.
\]
\end{definition}

\begin{remark*}[Interpretation]
The statement \(x > y\) means that \(x\) is obtained from \(y\) by adding a natural number.
\end{remark*}

\begin{definition}[Less than]
\label{def:less-than}
Let \(x,y \in \mathbb{N}\). We write
\[
x < y
\]
if there exists \(v \in \mathbb{N}\) such that
\[
y = x + v.
\]
\end{definition}

\begin{remark*}[Interpretation]
The statement \(x < y\) means that \(y\) is obtained from \(x\) by adding a natural number.
\end{remark*}

\begin{theorem}[Trichotomy of Order]
\label{thm:trichotomy-of-order}
For any \(x,y \in \mathbb{N}\), exactly one of the following holds:
\[
x = y,
\qquad
x > y,
\qquad
x < y.
\]
\end{theorem}

\begin{remark*}[Interpretation]
Any two natural numbers are comparable in exactly one way: they are equal, or the first is greater, or the first is less.
\end{remark*}

\begin{theorem}[Symmetry of Strict Greater Than]
\label{thm:symmetry-of-greater-than}
Let \(x,y \in \mathbb{N}\). If
\[
x > y,
\]
then
\[
y < x.
\]
\end{theorem}

\begin{remark*}[Interpretation]
Reversing a strict inequality reverses the order relation.
\end{remark*}

\begin{theorem}[Symmetry of Strict Less Than]
\label{thm:symmetry-of-less-than}
Let \(x,y \in \mathbb{N}\). If
\[
x < y,
\]
then
\[
y > x.
\]
\end{theorem}

\begin{remark*}[Interpretation]
The relation \(<\) is the same comparison as \(>\), read in the opposite direction.
\end{remark*}

\begin{definition}[Greater than or equal to]
\label{def:greater-than-or-equal-to}
Let \(x,y \in \mathbb{N}\). We write
\[
x \ge y
\]
if
\[
x > y
\quad \text{or} \quad
x = y.
\]
\end{definition}

\begin{remark*}[Interpretation]
The relation \(x \ge y\) means that \(x\) is either strictly greater than \(y\) or equal to \(y\).
\end{remark*}

\begin{definition}[Less than or equal to]
\label{def:less-than-or-equal-to}
Let \(x,y \in \mathbb{N}\). We write
\[
x \le y
\]
if
\[
x < y
\quad \text{or} \quad
x = y.
\]
\end{definition}

\begin{remark*}[Interpretation]
The relation \(x \le y\) means that \(x\) is either strictly less than \(y\) or equal to \(y\).
\end{remark*}

\begin{theorem}[Symmetry of Weak Greater Than]
\label{thm:symmetry-of-greater-than-or-equal-to}
Let \(x,y \in \mathbb{N}\). If
\[
x \ge y,
\]
then
\[
y \le x.
\]
\end{theorem}

\begin{remark*}[Interpretation]
Reversing a weak inequality exchanges \(\ge\) and \(\le\).
\end{remark*}

\begin{theorem}[Symmetry of Weak Less Than]
\label{thm:symmetry-of-less-than-or-equal-to}
Let \(x,y \in \mathbb{N}\). If
\[
x \le y,
\]
then
\[
y \ge x.
\]
\end{theorem}

\begin{remark*}[Interpretation]
The relation \(\le\) is the reverse reading of \(\ge\).
\end{remark*}

\begin{theorem}[Transitivity of Strict Ordering]
\label{thm:transitivity-of-strict-ordering}
Let \(x,y,z \in \mathbb{N}\). If
\[
x < y
\quad \text{and} \quad
y < z,
\]
then
\[
x < z.
\]
\end{theorem}

\begin{remark*}[Interpretation]
If one natural number is below a second, and the second is below a third, then the first is below the third.
\end{remark*}

\begin{theorem}[Mixed Transitivity]
\label{thm:mixed-transitivity}
Let \(x,y,z \in \mathbb{N}\). If either
\[
x \le y
\quad \text{and} \quad
y < z,
\]
or
\[
x < y
\quad \text{and} \quad
y \le z,
\]
then
\[
x < z.
\]
\end{theorem}

\begin{remark*}[Interpretation]
A strict inequality is preserved if the other comparison in the chain is weak.
\end{remark*}

\begin{theorem}[Transitivity of Weak Ordering]
\label{thm:transitivity-of-weak-ordering}
Let \(x,y,z \in \mathbb{N}\). If
\[
x \le y
\quad \text{and} \quad
y \le z,
\]
then
\[
x \le z.
\]
\end{theorem}

\begin{remark*}[Interpretation]
Weak order is transitive.
\end{remark*}

\begin{theorem}[Addition Produces a Greater Number]
\label{thm:addition-produces-a-greater-number}
Let \(x,y \in \mathbb{N}\). Then
\[
x + y > x.
\]
\end{theorem}

\begin{remark*}[Interpretation]
Adding a natural number to \(x\) produces a number strictly greater than \(x\).
\end{remark*}

\begin{theorem}[Addition Preserves Order]
\label{thm:addition-preserves-order}
Let \(x,y,z \in \mathbb{N}\).
\begin{enumerate}[label=(\alph*)]
\item If
\[
x > y,
\]
then
\[
x + z > y + z.
\]

\item If
\[
x = y,
\]
then
\[
x + z = y + z.
\]

\item If
\[
x < y,
\]
then
\[
x + z < y + z.
\]
\end{enumerate}
\end{theorem}

\begin{remark*}[Interpretation]
Adding the same natural number to both sides preserves the comparison.
\end{remark*}